\documentclass[10pt,a4paper]{amsart}
\usepackage{amsmath,amssymb,amsthm,graphicx}
% Retained minimal preamble for the Stacks Project chapter more-algebra.tex (More on Algebra).
% The chapter file pinned-original-source.tex carries only \input{preamble} on its first line, so no
% preamble text lives inside the chapter itself.  The Stacks Project uses one single shared preamble for
% every chapter; the snapshot stages it identically as sources/stacks-chow/preamble.tex and
% sources/stacks-etale-cohomology/preamble.tex.  The block below is copied mechanically, line for line,
% from sources/stacks-chow/preamble.tex lines 14-18 (the Xy-pic setup) and 154-217 (the native theorem
% declarations together with the mathematical macros).  Nothing was added, removed or redefined; the
% wrapper class and the amsmath/amssymb/amsthm/graphicx base packages are supplied by the assembler.
% Chapter: more-algebra.tex -- More on Algebra.
% For commutative diagrams we use Xy-pic
\usepackage[all]{xy}

% We use 2cell for 2-commutative diagrams.
\xyoption{2cell}

% Theorem environments.
%
\theoremstyle{plain}
\newtheorem{theorem}[subsection]{Theorem}
\newtheorem{proposition}[subsection]{Proposition}
\newtheorem{lemma}[subsection]{Lemma}

\theoremstyle{definition}
\newtheorem{definition}[subsection]{Definition}
\newtheorem{example}[subsection]{Example}
\newtheorem{exercise}[subsection]{Exercise}
\newtheorem{situation}[subsection]{Situation}

\theoremstyle{remark}
\newtheorem{remark}[subsection]{Remark}
\newtheorem{remarks}[subsection]{Remarks}

\numberwithin{equation}{subsection}

% Macros
%
\def\lim{\mathop{\mathrm{lim}}\nolimits}
\def\colim{\mathop{\mathrm{colim}}\nolimits}
\def\Spec{\mathop{\mathrm{Spec}}}
\def\Hom{\mathop{\mathrm{Hom}}\nolimits}
\def\Ext{\mathop{\mathrm{Ext}}\nolimits}
\def\SheafHom{\mathop{\mathcal{H}\!\mathit{om}}\nolimits}
\def\SheafExt{\mathop{\mathcal{E}\!\mathit{xt}}\nolimits}
\def\Sch{\mathit{Sch}}
\def\Mor{\mathop{\mathrm{Mor}}\nolimits}
\def\Ob{\mathop{\mathrm{Ob}}\nolimits}
\def\Sh{\mathop{\mathit{Sh}}\nolimits}
\def\NL{\mathop{N\!L}\nolimits}
\def\CH{\mathop{\mathrm{CH}}\nolimits}
\def\proetale{{pro\text{-}\acute{e}tale}}
\def\etale{{\acute{e}tale}}
\def\QCoh{\mathit{QCoh}}
\def\Ker{\mathop{\mathrm{Ker}}}
\def\Im{\mathop{\mathrm{Im}}}
\def\Coker{\mathop{\mathrm{Coker}}}
\def\Coim{\mathop{\mathrm{Coim}}}

% Boxtimes
%
\DeclareMathSymbol{\boxtimes}{\mathbin}{AMSa}{"02}

%
% Macros for moduli stacks/spaces
%
\def\QCohstack{\mathcal{QC}\!\mathit{oh}}
\def\Cohstack{\mathcal{C}\!\mathit{oh}}
\def\Spacesstack{\mathcal{S}\!\mathit{paces}}
\def\Quotfunctor{\mathrm{Quot}}
\def\Hilbfunctor{\mathrm{Hilb}}
\def\Curvesstack{\mathcal{C}\!\mathit{urves}}
\def\Polarizedstack{\mathcal{P}\!\mathit{olarized}}
\def\Complexesstack{\mathcal{C}\!\mathit{omplexes}}
% \Pic is the operator that assigns to X its picard group, usage \Pic(X)
% \Picardstack_{X/B} denotes the Picard stack of X over B
% \Picardfunctor_{X/B} denotes the Picard functor of X over B
\def\Pic{\mathop{\mathrm{Pic}}\nolimits}
\def\Picardstack{\mathcal{P}\!\mathit{ic}}
\def\Picardfunctor{\mathrm{Pic}}
\def\Deformationcategory{\mathcal{D}\!\mathit{ef}}
\usepackage{cleveref}
\crefname{originallemma}{original lemma}{original lemmas}
\Crefname{originallemma}{Original lemma}{Original lemmas}
\makeatletter
\newlabel{more-algebra-lemma-absolutely-integrally-closed-henselian-pair}{{0DCT}{source}{}{sourcepointer-1}{}}
\newlabel{more-algebra-lemma-absolutely-integrally-closed-henselian-pair@cref}{{[originallemma][0][]0DCT}{[1][0][]source}{}{}{}}
\makeatother
\title{An ideal in an absolutely integrally closed ring gives a henselian pair under two conditions}
\author{The Stacks Project Authors}
\date{stacks-project:more-algebra.tex:local-snapshot-20261012; GNU FDL 1.2 or later}
\begin{document}
\maketitle
\setcounter{subsection}{1}
\noindent\textit{Editorial note.} All delivered mathematical text is the author's own, copied from the pinned Stacks Project chapter file more-algebra.tex; no statement or proof step was rewritten, reordered or removed.\par\medskip
\noindent\textit{Editorial note.} The chapter file pinned-original-source.tex carries only \textbackslash{}input\{preamble\} on its first line, so no preamble text lives inside the chapter. The retained file sources/stacks-more-algebra/retained-preamble.tex is a mechanical, line-for-line copy of the project's single shared preamble as staged in the snapshot, namely sources/stacks-chow/preamble.tex lines 14-18 (Xy-pic setup) and 154-217 (theorem declarations and mathematical macros). Nothing was added, removed or redefined, no package was added, and the excerpt serves only as preamble, never as a span coordinate.\par\medskip
\noindent\textit{Editorial note.} Both counted spans start at the \textbackslash{}begin line of their environment, so the statement and the author's proof are delivered whole; surrounding prose of the section is deliberately outside the counted spans and is carried as standard prerequisites instead.\par\medskip
\noindent\textit{Editorial note.} The citation fields (source\_id, source\_version, authors, source\_url, license.name, source\_title) are identical to those of the cards already published from this chapter, so that the batch selector treats them as one source instead of deferring them as a same-source clash.\par\medskip
\noindent\textit{Allowed standard prerequisites (editorial).}\begin{itemize}
\item Absolutely integrally closed ring, meaning every monic polynomial decomposes into linear factors. definition environment at source lines 3509-3515 of the same chapter. Supplies the roots a\_i of f in A, which are the points at which the idempotents are evaluated.
\item Henselian pair, defined by the lifting of factorisations of monic polynomials modulo the ideal. the definition used by the chapter, cited inside the surrounding text of source lines 3580-3618. Fixes the goal: any factorisation modulo I with monic coprime factors must lift.
\item Idempotents of a ring correspond to decompositions of the ring as a finite product of rings. used at proof lines 3640-3647 where A is written as a finite product with each e\_i equal to 0 or 1. Reduces the problem to the case where each evaluated idempotent is 0 or 1, which is what makes the induction on the degree possible.
\end{itemize}\medskip
\ifdefined\hypertarget\hypertarget{sourcepointer-1}{}\fi
\noindent\textit{Source reference.} Original lemma 0DCT: focused-seeds/AGP/sources/stacks-snapshot-20261011/tags.txt, source lines 16743--16743. The lemma's own canonical tag, quoted so that readers can resolve the label inside the delivered text; it points at the counted statement itself and carries no proof content.\par\medskip
\begin{lemma}
\label{lemma-absolutely-integrally-closed-henselian-pair}
Let $A$ be absolutely integrally closed. Let $I \subset A$ be an ideal.
Then $(A, I)$ is a henselian pair if (and only if) the following
conditions hold
\begin{enumerate}
\item $I$ is contained in the Jacobson radical of $A$,
\item $A \to A/I$ induces a bijection on idempotents.
\end{enumerate}
\end{lemma}

\begin{proof}
Let $f \in A[T]$ be a monic polynomial and let
$f \bmod I = g_0 h_0$ be a factorization over $A/I$
with $g_0$, $h_0$ monic such that
$g_0$ and $h_0$ generate the unit ideal of $A/I[T]$.
This means that
$$
A/I[T]/(f) = A/I[T]/(g_0) \times A/I[T]/(h_0)
$$
Denote $e \in A/I[T]/(f)$ the element correspoing to the
idempotent $(1, 0)$ in the ring on the right.
Write $f = (T - a_1) \ldots (T - a_d)$ with $a_i \in A$.
For each $i \in \{1, \ldots, d\}$ we obtain an $A$-algebra map
$\varphi_i : A[T]/(f) \to A$, $T \mapsto a_i$ which induces
a similar $A/I$-algebra map $\overline{\varphi}_i : A/I[T]/(f) \to A/I$.
Denote $e_i = \overline{\varphi}_i(e) \in A/I$. These are idempotents.
By our assumption (2) we can lift $e_i$ to an idempotent in $A$.
This means we can write $A = \prod A_j$ as a finite product of
rings such that in $A_j/IA_j$ each $e_i$ is either $0$ or $1$.
Some details omitted. Observe that $A_j$ is absolutely integrally
closed as a factor ring of $A$. It suffices to lift the factorization
of $f$ over $A_j/IA_j$ to $A_j$.
This reduces us to the situation discussed in the next paragraph.

\medskip\noindent
Assume $e_i = 1$ for $i = 1, \ldots, r$ and $e_i = 0$ for
$i = r + 1, \ldots, d$. From $(g_0, h_0) = A/I[T]$ we have that there are
$k_0, l_0 \in A/I[T]$ such that $g_0 k_0 + h_0 l_0 = 1$.
We see that $e = h_0 l_0$ and $e_i = h_0(a_i) l_0(a_i)$.
We conclude that $h_0(a_i)$ is a unit for $i = 1, \dots ,r$.
Since $f(a_i) = 0$ we find $0 = h_0(a_i)g_0(a_i)$
and we conclude that $g_0(a_i) = 0$ for $i = 1, \ldots, r$.
Thus $(T - a_1)$ divides $g_0$ in $A/I[T]$, say
$g_0 = (T - a_1) g_0'$.
Set $f' = (T - a_2) \ldots (T - a_d)$ and $h'_0 = h_0$.
By induction on $d$ we can lift the factorization $f' \bmod I = g'_0 h'_0$
to a factorization of $f' = g' h'$ over over $A$ which
gives the factorization $f = (T - a_1) g' h'$
lifting the factorization $f \bmod I = g_0 h_0$
as desired.
\end{proof}

\end{document}
